\section{A priori error estimates}
\label{sec:error-estimates}

We first isolate the abstract stability argument.

\begin{theorem}[C\'ea's lemma]
\label{thm:cea}
Let $u\in\V$ and $u_h\in\Vh$ solve the continuous and discrete problems.
Under \cref{eq:continuity,eq:coercivity},
\begin{equation}
  \norm{u-u_h}_{\V}
  \leq \frac{M}{m}
  \inf_{v_h\in\Vh}\norm{u-v_h}_{\V}.
  \label{eq:cea-estimate}
\end{equation}
\end{theorem}

\begin{proof}
For arbitrary $v_h\in\Vh$, coercivity and Galerkin orthogonality imply
\begin{align*}
  m\norm{u-u_h}_{\V}^2
  &\leq a(u-u_h,u-u_h)\\
  &=a(u-u_h,u-v_h)\\
  &\leq M\norm{u-u_h}_{\V}\norm{u-v_h}_{\V}.
\end{align*}
Cancel the error norm and take the infimum over $v_h$.
\end{proof}

Let $I_hu$ denote the nodal interpolant. For $u\in H^2(\Omega)$ and a
shape-regular mesh,
\begin{equation}
  \norm{u-I_hu}_{H^1(\Omega)}
  \leq C_{\mathrm{int}}h\seminorm{u}_{H^2(\Omega)}.
  \label{eq:interpolation-estimate}
\end{equation}
Combining \cref{thm:cea,eq:interpolation-estimate} gives the principal energy
estimate.

\begin{corollary}[First-order energy convergence]
\label{cor:energy-rate}
If $u\in H^2(\Omega)\cap H_0^1(\Omega)$, then
\begin{equation}
  \norm{u-u_h}_{H^1(\Omega)}
  \leq C h\seminorm{u}_{H^2(\Omega)},
  \label{eq:h1-error}
\end{equation}
where $C$ is independent of $h$.
\end{corollary}

An additional regularity assumption produces a higher-order $L^2$ estimate.
For $e=u-u_h$, let $z\in\V$ solve the adjoint problem
\begin{equation}
  a(v,z)=\inner{e}{v}_{L^2(\Omega)}
  \qquad\text{for every }v\in\V.
  \label{eq:dual-problem}
\end{equation}
Assume elliptic regularity $\norm{z}_{H^2}\leq C_{\mathrm{reg}}\norm{e}_{L^2}$.
Using Galerkin orthogonality with $I_hz$ gives
\begin{align}
  \norm{e}_{L^2}^2
  &=a(e,z-I_hz)\notag\\
  &\leq C h\norm{e}_{H^1}\norm{z}_{H^2}
  \leq C h\norm{e}_{H^1}\norm{e}_{L^2}.
  \label{eq:duality-chain}
\end{align}
Together with \cref{eq:h1-error}, this proves
\begin{equation}
  \norm{u-u_h}_{L^2(\Omega)}
  \leq C h^2\seminorm{u}_{H^2(\Omega)}.
  \label{eq:l2-error}
\end{equation}

\begin{remark}
The $L^2$ rate depends on adjoint regularity and may deteriorate on nonconvex
domains with re-entrant corners. The energy estimate remains quasi-optimal
without this extra assumption.
\end{remark}
